\section{Chapter 7 --- Analysis: Mechanisms, Noise Floors, Failure Modes, and Ceilings}\label{sec:chapter-7-analysis-mechanisms-noise-floo}

Chapter 6 reported the scoreboard; this chapter explains the mechanisms behind it. The aim is that the reader believes the results not because the deltas are large, but because the failure structure beneath them is accounted for. Seven analyses carry that burden. Section 7.1 gives a mechanism-explained negative result for preference optimisation. Section 7.2 traces the four-point dev-to-test decay pattern, which identifies distillation-as-denoising as the operative mechanism, and adds a post-hoc cross pairing that isolates briefing quality as the dominant variable. Section 7.3 measures the teacher provider-noise floor that disciplines every teacher-relative claim in this dissertation. Section 7.4 decomposes the error modes and identifies a structural ceiling the students share with their teacher. Section 7.5 reports a schema-format diagnostic of what fine-tuning actually teaches, together with its own instrumentation history as a methodological finding. Section 7.6 examines the hard-composite slice on which the student's advantage is largest. Section 7.7 discloses the pilot results of a pre-registered compositional probe, exactly as far as they reached.

\section{DPO and the toxicity of near-miss negatives}\label{sec:dpo-and-the-toxicity-of-near-miss-negati}

Why did DPO help on one base and hurt on the other? At the close of Chapter 6, Direct Preference Optimisation was the champion rung of the Qwen ladder: DPO-v1 reached 83.5\% on the development set (217/260), a net +2.0 points over the RSFT rung beneath it (81.5\%). That aggregate hides a regression which the per-bucket decomposition exposes at once. On the hardest question family, bridge\_join, the RSFT rung answered 16 of 20 correctly and the DPO rung only 11. The preference stage gained in aggregate while losing on exactly the capability the preference pairs were richest in. On the second base the same treatment was not just uneven but net-destructive: Llama DPO-v1 scored 77.3\% against its RSFT parent's 82.7\%, a --5.4 point regression. Read together, the two bases trace a dose--response curve for a single pathology --- mild on Qwen (+2.0 net, with a localised bridge regression) and severe on Llama (--5.4 net) --- rather than two unrelated tuning accidents.

The training internals of the Llama run identify the pathology. Its DPO loss saturated at 0.03 with pair-classification accuracy of 99.4\%. The log-probability of the chosen responses sat at --7.7, essentially memorised from the RSFT stage that preceded it. The growth in preference margins came entirely from pushing the log-probability of \emph{rejected} responses down, from --185 to --203. Nothing was being learned about the chosen plans; they were already at ceiling. The optimiser therefore satisfied the objective by displacing probability mass away from the rejected plans and, as we argue below, away from correct plans adjacent to them. The damage had a behavioural signature visible before any accuracy number was computed: during its evaluation the model's generation throughput collapsed by roughly a factor of nine, the symptom of output distributions flattened under guided decoding.

The natural response is to blame the recipe, and the recipe was genuinely open to criticism. The v1 pair pool had never been bucket-capped: 71\% of its pairs came from the bridge and set families, and 766 of its 1,079 pairs originated from oracle-gated repair traces. The loss also carried no anchor on the chosen responses (no auxiliary likelihood term). We therefore ran a two-arm recipe-fix ablation following contemporary curation practice. The pools were rebuilt on-policy-first (for Qwen, 389 pairs of which 94\% were on-policy), bucket-capped, and deduplicated to a balanced composition. Arm A used the sigmoid loss with an auxiliary chosen-likelihood term; arm B replaced the objective with IPO. Both trained for three epochs after a one-epoch pilot showed initial margins of --17: on-policy pairs are hard against a base-model reference, and the first run had simply been undertrained at 25 steps.

Both arms failed. Arm A scored 80.0\% with a bridge sub-score of 5/20; arm B scored 82.3\% with bridge 9/20. Both lose to the unfixed champion DPO-v1 (83.5\%, bridge 11/20), and all three lose the bridge bucket to plain RSFT (16/20). The recipe fixes did not just fail to help. The \emph{cleaner}, more on-policy, better-balanced pools produced \emph{worse} bridge performance than the naive pool they replaced.

That inversion is the diagnostic fact, and it points to a mechanism rather than a tuning error. The rejected bridge plans in an on-policy pool are \emph{near-miss} negatives. They are string-adjacent to correct bridge plans: the same four-variable dataflow, differing only in the return target or a single slot. When the objective suppresses them, it also drags down the probability of the correct modes next to them. An anchor term or the IPO objective cannot protect what it cannot separate, because the protected and suppressed manifolds coincide. This is consistent with the likelihood-displacement analysis of Razin et al., in which preference optimisation is most destructive exactly when chosen and rejected responses are similar. On this reading, balancing and on-policy filtering \emph{concentrated} the poison: they raised the fraction of pairs whose negative sat closest to the correct answer. The dose--response across bases also follows naturally. The Llama policy, whose RSFT stage had already memorised the chosen set, had nothing left to gain and everything to displace.

Two consequences were adopted, and both are visible elsewhere in this dissertation. First, the claim is bounded. These experiments do not show that DPO is toxic in general; they show it is hazardous in \emph{near-miss-negative regimes}, where the verifier's rejects are structurally adjacent to its passes --- which is exactly the regime a verifier-gated pipeline produces. Second, the bridge-capability lever was moved from stronger suppression to \emph{additive} learning: the r2 bootstrap round (Chapter 6) consolidates correct bridge plans by RSFT-style self-distillation from the DPO policy, with no negative term at all. Within the terms of Pillar 3, the finding is a controlled contrast between subtractive and additive use of the same verifier signal. Filtering what enters the positives is safe; suppressing what the filter rejects is not, whenever the rejects live next door to the truth.

\section{The four-point decay pattern: distillation as denoising}\label{sec:the-four-point-decay-pattern-distillatio}

Why do the hybrid systems decay from the development set to the test set while the fully-local systems stay flat? The most consequential single pattern in the confirmatory results is not any one accuracy but the \emph{shape} of this decay across the four student systems. Both fully-local systems are essentially flat from the development set to the held-out test set: Qwen fully-local decays --0.5 points (86.2\% $\rightarrow$ 85.65\%) and Llama fully-local --0.9 points (84.2\% $\rightarrow$ 83.33\%). Both hybrid systems fall steeply on the identical transition: Qwen hybrid --5.5 points (83.5\% $\rightarrow$ 78.03\%) and Llama hybrid --7.1 points (83.1\% $\rightarrow$ 75.97\%). The pattern is four-for-four and cross-base: flat wherever the understanding stage (Step 1) is the distilled local adapter, steep wherever it is the cloud model, with the \emph{same} planning checkpoints on each base. Whatever explains the decay must therefore live in the briefing stage, not in the planner and not in the base model.

Three strands of evidence, each with its own experiment, converge on one account.[\^{}strands] First, composition. The final test set's natural question mix upweights the bridge and comparison families to 28.7\% of items, against 15.4\% under the development set's uniform bucketing. These are exactly the families where the structure and regularity of the intermediate briefing most strongly determine whether the planner can produce a correct multi-variable plan. Second, distributional match. The cloud model's briefings on hard questions are \emph{irregular} --- free-form in structure and emphasis --- whereas the student's Step-2 planner was trained on the regular briefing distribution. The distilled Step-1's output lands inside the distribution the planner knows; the cloud model's improvisations land outside it. Third, boundary behaviour. An autopsy of the hybrid system's answerable-to-null down-flips on the test set found that four of six were \emph{deterministically} reproducible conservative abstentions (boundary questions pushed to the abstain side by irregular briefings), with only two of six attributable to provider churn. The three strands are one mechanism seen at three grains: the harder mix concentrates weight on briefing-sensitive questions, irregular briefings fall outside the planner's trained distribution there, and at the decision boundary that mismatch resolves into abstention. Distillation of the understanding stage acts as \emph{denoising}: the Step-1 student, trained only on briefings whose entire downstream pipeline outcome was verified and oracle-correct, reproduces the teacher's competence stripped of its irregularity.[\^{}fileorder]

[\^{}strands]: Strand one is a distributional census of final\_test versus the development buckets; strand two rests on the Step-1 distillation filter described in Chapter 6 (5,091 briefings retained solely on whole-pipeline verified-and-oracle-correct outcome); strand three is the down-flip autopsy of 2026-07-07, in which the six answerable$\rightarrow$null flips were re-run through a live hybrid configuration: 4/6 reproducibly null, 2/6 flipped back correct.

[\^{}fileorder]: A rolling-window audit initially suggested a mid-run incident: both the hybrid system and the teacher show their worst accuracy window near file position 850. The audit resolved this benignly. Positions 800--1000 of final\_test are a 200/200 cluster of bridge\_join items in file order, so the co-located dip is that cluster, not a temporal API event; the zero-cloud fully-local run swings at the same positions.

One further observation sharpens the conclusion, and it must be read with its label attached: it is a post-hoc pairing, not a pre-registered comparison.[\^{}posthoc] The weaker base with the local briefing beats the stronger base with the cloud briefing. Llama fully-local scores 83.33\% against Qwen hybrid's 78.03\% (discordant pairs +252/--131, $\Delta$ = +5.30 points, 95\% CI [+3.62, +6.97], $p=8.7\times10^{-10}$), even though Qwen sits 2--3 points above Llama in every same-briefing configuration in the matrix. Briefing quality is the dominant variable: the quality difference between the two Step-1 stages outweighs the capability difference between the two base models.

[\^{}posthoc]: This cross pairing was noticed after the five-pairing matrix completed and was not pre-registered; it is reported as a mechanism observation only and carries no headline weight.

\section{The teacher noise floor}\label{sec:the-teacher-noise-floor}

How much does the teacher differ from itself? Every claim in this dissertation of the form "system X differs from the teacher by $\Delta$" presupposes an answer to that question. Three identical v2.2 evaluations of the teacher pipeline on the development set --- same questions, same configuration, different nights --- returned 71.9\%, 71.9\%, and 73.9\%: a mean of 72.6\% with a standard deviation of 1.1 points. The stability of the totals is deceptive. Pairwise, the three replicates disagree on between 13 and 22 individual questions. The provider's unseeded temperature-zero decoding churns roughly 5--8\% of predictions between runs that are configured identically, and the churn largely cancels in the aggregate.

Two disciplines follow, and both are enforced throughout. First, no single-run teacher delta below roughly three points is claimable. Teacher-relative headline claims either clear the floor by a wide margin (the student--teacher gaps of Chapter 6 exceed it several times over) or are paired against \emph{each} replicate separately. Second, an honest account of where we were caught by this ourselves. During the v2.0$\rightarrow$v2.1 pipeline iteration we chased an apparent teacher regression (70.0\% $\rightarrow$ 66.5\% $\rightarrow$ 66.1\%) as if it were entirely a scaffolding defect. The later adversarial review established that the chase was partly noise-chasing. The v2.1 gate-narrowing was a real fix, with deterministic answerable-to-null signatures behind it, but the residual 66.1\% was a bad draw from the noise floor, not evidence of remaining gate damage. We record the episode as a lesson because it is the mirror image of the discipline: attributing a within-floor fluctuation to a cause is as much an error as claiming one.

The floor also disposes of a tempting misreading of the confirmatory set. The teacher's 69.76\% on final\_test, below its 72.6\% development mean, is difficulty, not API degradation. An integrity audit compared the teacher's predictions on the 260 development questions embedded within final\_test across the two runs and found the teacher \emph{net positive} (+6) on the later run. The run that produced 69.76\% was, if anything, a slightly favourable draw; the lower total reflects the harder natural mix of the full test set.

\section{Error-mode decomposition and the boolean structural ceiling}\label{sec:error-mode-decomposition-and-the-boolean}

Which errors does scaffolding fix, which does training fix, and what remains when both levers are exhausted? The v1 evaluation matrix supplied the error taxonomy from which the v2 pipeline was built. Of the champion rung's development-set errors, 21 were abstentions on answerable questions, 12 were answers on unsupported questions, and 22 were wrong values. Chapter 5's worked example traced these to four precise scaffolding defects: an over-strict type gate on empty-value lookups, an over-permissive alias substring rule, a consistency-check bypass on the compiled fast path, and a verifier-invisible count-of-intermediate-set misread. Each was fixed deterministically and regression-tested against real pipeline intermediates. The decomposition's purpose in this chapter is the division of labour it establishes for Pillar 2. The first two error modes are dominated by \emph{scaffolding}: they moved when the checks moved, for every system including the untrained zero-shot rung. The wrong-value mode is dominated by \emph{planning quality} and moved mainly with training. What scaffolding fixes, it fixes for all systems at once (the v2.2 zero-shot floor rose to 70.4\% from v1's 61.5\%); what training fixes, it fixes per rung, on top of whatever floor the scaffolding provides.

What remains when both levers are exhausted is visible in one bucket with unusual clarity. On the development set's boolean family (n=20), nearly every system in the matrix lands on 14/20: the fully-local Qwen champion, the fully-local Llama system, and the hybrid champion alike (Llama-SFT and one teacher replicate land on 13/20). More telling than the matching totals is the matching items. The \emph{same six questions} are missed by the fully-local Qwen system, the fully-local Llama system, and the hybrid champion, six of six overlapping. All six are also missed by every one of the three teacher replicates.[\^{}booleansrc] This is not a training gap. No rung of either ladder, at any training dose, touches these questions, and neither does the teacher under any draw. It is a shared structural ceiling: some property these six questions require is not currently expressible through the pipeline's plan language or its checks, for teacher and students alike. Distillation cannot exceed what neither the teacher nor the scaffolding can express. The ceiling is pipeline-bound, and raising it is scaffolding work (moving the projection boundary), not learning work. We state this as a shared expressiveness ceiling and no more. The evidence does not distinguish which component of the chain (plan vocabulary, grounding, or the questions' own construction) carries the limitation, only that it is common to every system evaluated.

[\^{}booleansrc]: Computed from the per-item results files of the v2.2 matrix (outputs/eval/matrix\_v2/*/compare\_cicada.results.jsonl).

\section{The schema-format diagnostic: what does fine-tuning actually teach?}\label{sec:the-schema-format-diagnostic-what-does-f}

Does fine-tuning teach planning, or only the output format? The training ladder of Chapter 6 measures accuracy under guided decoding, which enforces the plan schema at generation time. This invites a deflationary reading: perhaps fine-tuning teaches nothing but the format, and the ladder's gains are an artefact of format acquisition. The schema diagnostic was designed to test this, with success criteria pre-committed before any output was read. Its result is a clean separation, but the path to that result matters as much as the result, and we report both.

The diagnostic evaluates three arms --- the untrained Qwen3-8B base, the SFT rung, and the DPO rung --- on 100 development questions with \emph{no} guided decoding. Each output is scored on two metrics: \emph{grounded-any-schema} (does the output contain a question-grounded plan in any schema at all?) and \emph{target-contract-conformant} (does it match the exact contract the compiler consumes?). The result: grounded-any is 98/98/99 across base/SFT/DPO, but target-contract conformance is 3/98/99. Opening ten randomly chosen base-arm outputs (a requirement, not a courtesy; see below) confirmed what the numbers suggest. The untrained base emits question-grounded plans with \emph{real} entities, such as Wessex Water, NHS England, and genuine CPV codes and years, organised into a coherent, self-invented schema (op:filter\_records, args.filters, id/inputs) that the compiler cannot consume. The base can plan; it cannot emit the executable contract. Fine-tuning's measured contribution is therefore \emph{conformance to the specific executable contract} (3 $\rightarrow$ 98/99), not planning ability itself. This is exactly the control-variable result the ladder needs. Guided decoding at evaluation time enforces the one thing that is near-free for trained models and near-impossible for the base, so the accuracy ladder isolates planning \emph{quality within the contract} rather than the ability to emit it. Two honesty caveats attach. The diagnostic stored only failures, so trained-arm passes were not individually inspected; the claim rests on the base arm's verified grounded-but-off-schema pattern. And one base abstention expressed through a return-operation field was under-counted as empty. The 3-versus-98 direction is insensitive to both.

The raw-inspection requirement exists because this diagnostic produced three consecutive wrong answers before it produced a right one, and each wrong answer was caught only by opening raw outputs. The first instrument \emph{over-counted} base-arm parse failures: a greedy JSON extractor was confused by markdown fences. We patched it. The second \emph{over-counted} base-arm successes: a well-formedness gate scored 100/100 for the base, and four spot-checked outputs revealed these to be identical degenerate shells, structurally present but semantically empty, so a content gate was added. The third \emph{under-counted} the base: the content gate looked for filters only at the trained schema's location, scored the base's genuinely grounded plans as zero, and produced a suspiciously clean 0/94/99. We retracted that result before publication and replaced it with the dual-metric design above. Three gate bugs, three opposite-signed biases, each one capable of carrying a false conclusion into this dissertation. Each was caught by the same act: refusing to cite a gate-produced number without reading the raw outputs behind it. We generalised the episode into a standing rule for every gate-decided number in this work --- ten deterministic-random raw outputs, passes and failures, human-read before the number enters the results table. We offer the audit trail itself as an evaluation-methodology contribution. Every gate is a measurement instrument, and an unverified instrument that produces a clean number is a liability precisely \emph{because} the number is clean.

\section{The hard-composite slice: where the student's advantage is largest}\label{sec:the-hard-composite-slice-where-the-stude}

Where is the student's advantage over the teacher largest? Chapter 6 established that the fully-local student beats the teacher on aggregate; the slice analysis locates where. The development-era analysis partitions the evaluation set into an iid slice (n=103) and a hard-composite slice (n=136). The latter comprises the hard operator families, second-level paraphrase rewrites, and abstention classes. One naming discipline governs everything said about it: this is a \emph{difficulty composite}, not a held-out distribution. The class is present in both training and test (with disjoint plans and novel surfaces), so no claim of out-of-distribution generalisation is made or licensed here; compositional novelty is the business of the pre-registered probe in §7.7.

\begin{table}[t]
\centering
\small
\caption{Accuracy on the iid versus hard-composite slices of the development set (absolute values; the pattern, not single points, carries the claim).}
\label{tab:07_analysis_1}
\begin{tabular}{llll}
\toprule
System & iid (n=103) & hard-composite (n=136) & slope \\
\midrule
Teacher r1 / r2 / r3 & 75.7\% / 73.8\% / 78.6\% & 65.4\% / 67.6\% / 66.2\% & --10.3 / --6.1 / --12.5 \\
Qwen DPO-v1 (hybrid) & 86.4\% & 78.7\% & --7.7 \\
Qwen fully-local & 84.5\% & 85.3\% & +0.8 \\
\bottomrule
\end{tabular}
\end{table}


The pattern, not any single cell, is the finding. Every teacher replicate decays by 6 to 12 points from the iid slice to the hard composite; the hybrid student decays by 7.7; the fully-local student shows no measurable decay (+0.8 points, n=136). The fully-local system uses the \emph{same} planning checkpoint as the hybrid, so the difference is attributable to the distilled understanding stage. On the hard-composite slice itself the fully-local student's advantage over the teacher reaches 18--20 points, its largest margin anywhere in the evaluation. The student trained only on verifier-passed supervision holds its accuracy exactly where the far larger teacher loses it. This is the mechanism evidence for the one-sentence thesis in its sharpest available form. The caveats are stated with the finding: slice sizes are modest (103 and 136), the slopes carry corresponding uncertainty, and the wording "no measurable decay" describes this pattern at this n; it is not a claim of proven invariance. The four-point dev-to-test pattern of §7.2 --- the same flat-versus-steep signature at n=2,285 --- is the larger-scale corroboration that keeps this small-n table from bearing weight alone.

\section{The compositional probe: pilot findings and a boundary characterisation}\label{sec:the-compositional-probe-pilot-findings-a}

Does the system generalise to plan compositions it has never seen? The claims above stop deliberately short of that question, which was assigned to a separate pre-registered instrument: ood\_probe\_v1. The probe has five question templates whose structural signatures --- rebuilt from gold-plan flat specifications, not template names --- are absent from the entire corpus: three bridge compositions (B1 bridge\_argmax, B2 bridge\_top\_k, B3 bridge\_compare) and two non-bridge compositions (N1 grouped temporal argmax, N2 per-side filtered sum-compare), with three-branch success criteria committed before generation. The probe was \emph{paused at the pilot stage} by an explicit scope decision before full generation began. No part of the planned \textasciitilde{}600-row evaluation was run, and nothing in this dissertation cites it as a generalisation result. The pilot itself, however, produced findings worth reporting exactly as far as they go.

First, the novelty gate passed its self-audit 10/10 in both directions: five existing question families were correctly recognised as present in the corpus signature inventory, and all five probe templates as absent. Second, the pilot corrected the pre-registration itself. N2's claimed novelty on the comparison-side-type axis proved false on faithful reconstruction: its comparison side is an existing type, and the template is novel only on the filter-slot axis, which makes it materially the weakest of the five. The correction is recorded against the pre-registration rather than papered over. (The reconstruction also revised an ad-hoc signature census from 50 to 38 under the committed five-field definition; novelty assertions are unaffected under both granularities.)

Third, and scientifically the most interesting, the executability pilot ran five hand-written probe questions through the exact production entry point of the fully-local pipeline. That mandate has a history: an earlier hand-wired bypass of the production path had produced a nearly \emph{inverted} executability table and was voided. The lesson, of a piece with §7.5, is that a pilot that does not enter through the front door measures a different system. Through the production path, all three bridge compositions executed end-to-end, including the highest-risk B3, in which a comparison decomposes into two bridged sub-plans. Both non-bridge compositions failed, for genuine and informative compile-time reasons. The executor's argmax returns an extremum record and offers no grouped-argmax, so N1 is rejected at intent compilation. The two-sided sum-comparison of N2 drops the per-side year filters that distinguish its arms. The system's compositional coverage is therefore \emph{asymmetric}. Novel compositions within the bridge family --- the family the training signal was richest in --- execute, while the two non-bridge aggregation compositions chosen for the probe fall outside the executor's current expressible plan space. We report this as a boundary characterisation of the executor, not as a generalisation result in either direction. It says where a full probe could measure generalisation, and where the pipeline itself, not the learner, is the binding constraint: a compositional cousin of the structural ceiling of §7.4.

\section{Synthesis}\label{sec:synthesis}

The analyses compose into a single picture. The verifier-gated pipeline creates an unusual learning environment. Its positives are clean, which is why additive distillation keeps working (§7.2, §7.6). Its negatives are near-misses, which is why subtractive preference learning is hazardous (§7.1). Its teacher is noisy in a measurable, disciplinable way (§7.3). Its residual failures split cleanly into what scaffolding fixes, what training fixes, and what neither currently can (§7.4, §7.7). The diagnostic apparatus that produced these conclusions required the same scepticism as the systems it measured (§7.5). On this account, the headline numbers of Chapter 6 stand not as fortunate draws but as the visible surface of a mechanism --- denoised supervision from the verifiable region, learned competence extending past it. The failure modes of that mechanism have been located, measured, and where possible explained. Chapter 8 states the conclusions this evidence supports and the limits it does not cross.
